\SourceLeaf{203}\Needspace{10\baselineskip}\begin{center}
[4.]
\end{center}

[In \textsc{Lübsen}, \emph{Ausf\"uhrliches Lehrbuch der Elementar-Geometrie}, Hamburg 1851, page 11, we read the following.]

\{\textls{Explanation:} A straight line is one that does not change its position when it rotates about its two fixed endpoints\HistoricalNote{*)}{\{*)}{This is approximately how we once heard \textsc{Gauss} define a straight line while explaining the telescope and its proper use. This definition is fruitful in theory, as the propositions that immediately follow show. The stated criterion is also important in practice, e.\,g., in adjusting a telescope or boring a cylinder correctly, etc.\}}.\}



\Needspace{8\baselineskip}
\begin{center}
[5.]

\textnormal{THEORY OF THE PRESENTATION OF PROPOSITIONS CONCERNING SPATIAL RELATIONS.}
\end{center}

1) \textls{Draw} all the objects relevant to a proposition and designate them by letters or digits, points by one, lines by 2, surfaces by 3 and solids, if necessary, by 4, though one can often suffice here as well.

2) Draw separately the spaces that can move relative to one another when their motion enters the task discontinuously.

3) Designate the relevant relations of magnitude by letters as well.

4) Put what is to be proved in the form of an equation and work back step by step to its grounds.

The whole task is thus, as it were, broken down into its elements, which are then put together again in reverse order. One readily sees what can be condensed in the actual presentation.

